\documentclass[uplatex,dvipdfmx,a4paper,11pt]{jsarticle}
\usepackage[top=20mm,bottom=20mm,left=20mm,right=20mm]{geometry}
\usepackage{longtable,booktabs,array}
\usepackage{hyperref}
\usepackage{pxjahyper}
\usepackage{enumitem}
\usepackage{fancyhdr}
\usepackage{lastpage}
\setlength{\parindent}{0pt}
\setlength{\parskip}{0.55em}
\pagestyle{fancy}
\fancyhf{}
\lhead{第18章 工学基礎}
\rhead{SWEBOK v4.0a 日本語まとめ}
\cfoot{\thepage/\pageref{LastPage}}
\newcommand{\infobox}[2]{\vspace{0.4em}\noindent\fbox{\begin{minipage}{0.96\linewidth}\textbf{#1}\par\smallskip #2\end{minipage}}\vspace{0.4em}}
\newcommand{\todobox}[2]{\vspace{0.4em}\noindent\fbox{\begin{minipage}{0.96\linewidth}\textbf{TODO（図） #1}\par\smallskip\textbf{画像生成用プロンプト}\par #2\end{minipage}}\vspace{0.4em}}
\begin{document}
\begin{center}
{\Large\bfseries 第18章 工学基礎}\par
\vspace{0.4em}
Engineering Foundations の日本語まとめ\par
SWEBOK Guide v4.0a Chapter 18 を初学者向けに再構成\par
作成日：2026年5月6日\par
\end{center}
\hrule
\subsection*{この資料の主張}
ソフトウェア工学は、工学としての共通基盤を持つ。真の問題を見極め、制約の中で複数の解を比較し、証拠に基づいて判断し、結果を測定して改善することが重要である。第18章は、抽象化、実験、統計、測定、標準、根本原因分析を通じて、ソフトウェア技術者が工学的に判断するための土台を整理する。
\infobox{この資料の読み方}{専門用語は原則として日本語で示し、必要に応じて英語を括弧内に併記する。英語名は、元資料やツール上の名称を確認したいときの手がかりとして残す。初学者が前提知識なしに読めるように、各節では「何のための概念か」「実務では何を見ればよいか」を重視する。}
\tableofcontents
\clearpage
\section{全体概要：この章で学ぶこと}
この章は、ソフトウェア工学を単なるプログラミング技術ではなく、工学の一分野として理解するための基礎を扱う。工学では、現実の制約の中で問題を見極め、複数の解を考え、根拠を持って選び、実世界での結果を監視する。ソフトウェアでも同じである。要求、設計、実装、テスト、保守、運用のどの場面でも、抽象化、実験、統計、測定、標準、根本原因分析を使う必要がある。
\infobox{到達目標}{この章を読み終えると、工学プロセスの反復的な流れを説明できる。工学設計が一つの正解を探す活動ではなく、制約下でよい解を選ぶ活動であることを説明できる。抽象化、カプセル化、階層、代替的な抽象化を区別できる。経験的方法、統計分析、測定理論をソフトウェア工学に適用する理由を説明できる。標準、根本原因分析、インダストリー4.0とソフトウェア工学の関係を説明できる。}
\todobox{第18章の全体地図を入れる。}{白背景、モノクロ、A4本文幅に収まる横長の学習用図解を作成する。中央に「工学基礎」を置き、周囲に「工学プロセス」「工学設計」「抽象化とカプセル化」「経験的方法」「統計分析」「モデリング・シミュレーション・プロトタイピング」「測定」「標準」「根本原因分析」「インダストリー4.0」を配置する。各項目がソフトウェア開発の判断を支えることを矢印で示す。ラベルは日本語、細線、講義ノート風、余白を広めにする。}
\subsection{最初に必要な用語}
\begin{longtable}{>{\raggedright\arraybackslash}p{0.23\linewidth}>{\raggedright\arraybackslash}p{0.69\linewidth}}
\toprule
\textbf{用語} & \textbf{初学者向けの説明} \\
\midrule
\endfirsthead
\toprule
\textbf{用語} & \textbf{初学者向けの説明} \\
\midrule
\endhead
工学 & 現実世界の問題を、体系的・規律的・測定可能な方法で解く活動である。制約の中で有用な解を作る点が重要である。 \\
工学プロセス & 問題を理解し、基準を決め、候補解を挙げ、評価し、選び、実際の性能を監視する流れである。 \\
工学設計 & 望まれる必要や仕様を満たすシステム、構成要素、過程を、制約の中で考案する活動である。 \\
抽象化 & 細部を一時的に隠し、重要な概念や関係に注目すること。曖昧にすることではなく、扱う水準を明確にすることである。 \\
カプセル化 & 抽象化した対象の内部詳細を隠し、外部から使うための窓口だけを示す仕組みである。 \\
経験的方法 & 観察、実験、過去データの分析など、実際のデータに基づいて判断する方法である。 \\
統計分析 & 標本から母集団を推測し、ばらつきや関係を理解するための技法である。 \\
測定 & 属性を定めた方法で値や記号に対応させること。測った値をどの演算に使えるかは尺度に依存する。 \\
標準 & 要求、仕様、指針、用語、品質水準などを共有するための合意済みの基準である。 \\
根本原因分析 & 望ましくない結果の表面的な症状ではなく、再発を防ぐための深い原因を特定する方法群である。 \\
インダストリー4.0 & 製造業におけるデジタル化、AI、IoT、クラウド、ビッグデータなどを活用した新しい産業構造である。 \\
\bottomrule
\end{longtable}
\subsection{略語}
\begin{longtable}{>{\raggedright\arraybackslash}p{0.13\linewidth}>{\raggedright\arraybackslash}p{0.32\linewidth}>{\raggedright\arraybackslash}p{0.47\linewidth}}
\toprule
\textbf{略語} & \textbf{英語} & \textbf{日本語での意味} \\
\midrule
\endfirsthead
\toprule
\textbf{略語} & \textbf{英語} & \textbf{日本語での意味} \\
\midrule
\endhead
CAD & Computer-Aided Design & コンピュータ支援設計。設計図やモデルをコンピュータで作るための方法・道具である。 \\
CMMI & Capability Maturity Model Integration & 能力成熟度モデル統合。組織のプロセス成熟度を段階的に表す枠組みである。 \\
PDF & Probability Density Function & 確率密度関数。連続型の確率変数の分布を表す関数である。 \\
PMF & Probability Mass Function & 確率質量関数。離散型の確率変数が特定値を取る確率を表す関数である。 \\
RCA & Root Cause Analysis & 根本原因分析。再発防止のため、望ましくない結果の深い原因を調べる方法群である。 \\
SDLC & Software Development Life Cycle & ソフトウェア開発ライフサイクル。企画から保守までの開発活動の流れである。 \\
\bottomrule
\end{longtable}
\subsection{導入：なぜソフトウェア工学に工学基礎が必要か}
ソフトウェア工学は、コンピュータ科学の知識を使うだけでは成立しない。現実の問題を解くには、費用、納期、品質、安全性、標準、運用、社会的影響などを同時に扱う必要がある。これは他の工学分野と共通する。
第18章の目的は、他の知識領域を支える工学共通の考え方を整理することである。特に、問題解決としての工学プロセス、設計、抽象化、実験、統計、モデル、測定、標準、根本原因分析を、ソフトウェア工学で使える形で理解することに焦点がある。
\infobox{重要}{この章は、個別のプログラミング技術ではなく、ソフトウェア技術者が判断を誤らないための土台を扱う。工学基礎を理解すると、なぜ測るのか、なぜモデルを作るのか、なぜ標準に従うのか、なぜ根本原因まで調べるのかが見える。}
\section{1 工学プロセス}
工学プロセスは、すべての工学分野に共通する高水準の問題解決手順である。単純な一本道ではなく、途中で得た知識によって前の段階へ戻る反復的な流れである。見積りや仮説は誤ることがあるため、選んだ解の実世界での性能を監視し、必要なら別の案を再検討する。
\begin{longtable}{>{\raggedright\arraybackslash}p{0.26\linewidth}>{\raggedright\arraybackslash}p{0.66\linewidth}}
\toprule
\textbf{段階} & \textbf{考えること} \\
\midrule
\endfirsthead
\toprule
\textbf{段階} & \textbf{考えること} \\
\midrule
\endhead
真の問題を理解する & 依頼された問題が、実際に解くべき問題とは限らない。根本原因分析を使い、背後にある本当の問題を探す。 \\
選定基準を定義する & 費用、便益、納期、品質、安全性、標準、運用性など、候補解を比較する基準を明確にする。 \\
技術的に実現可能な候補解を挙げる & 最初に思いついた案だけで決めない。複数の実現可能な案を出し、よりよい案を見落とさない。 \\
候補解を基準に照らして評価する & 各案が、必要をどれだけ満たし、各基準をどれだけ満たすかを比較する。 \\
望ましい案を選ぶ & 基準に対して最もよい案を選択する。ただし、選択は推定に基づくため絶対ではない。 \\
選んだ案の性能を監視する & 実際に使った結果を観察し、期待と違えば原因を調べ、必要なら選択を見直す。 \\
\bottomrule
\end{longtable}
\todobox{工学プロセスの流れ図を入れる。}{白背景、モノクロ、A4本文幅の循環型フローチャートを作成する。「真の問題を理解する」から始まり、「選定基準を定義する」「実現可能な候補解を挙げる」「基準で評価する」「望ましい案を選ぶ」「性能を監視する」へ進む。最後から最初へ戻る反復矢印を描く。各箱に短い説明を添える。日本語ラベル、細線、講義ノート風。}
\infobox{ソフトウェアでの例}{障害対応では、再起動してサービスを戻すだけでは工学プロセスとして不十分である。真の問題がメモリリーク、設定ミス、負荷見積りの誤り、監視設計の不足のいずれかを調べ、複数の対策を比較し、実施後に再発率や復旧時間を監視する必要がある。}
\section{2 工学設計}
工学設計は、製品やシステムのライフサイクル費用を左右する。ソフトウェア設計でも、どの機能を実現するか、どの品質属性を達成するかによって、保守費用、障害対応費用、運用費用が大きく変わる。
工学の問題は、一つの正解を当てる問題ではないことが多い。多くの場合、複数の解があり、利用可能な資源、費用、標準、技術知識、物理法則、組織の制約などの中で、目的に最も合う実現可能な解を選ぶ。
\subsection{2.1 工学教育における工学設計}
工学教育の認定では、設計能力が重視される。学生は、複雑で開かれた問題に対して、システム、構成要素、過程を設計できることが求められる。その際、健康・安全、標準、経済、環境、文化、社会的考慮を無視できない。
この考え方はソフトウェアにも当てはまる。たとえば医療、金融、公共、交通、製造のソフトウェアでは、機能が動くだけでは足りない。標準、規制、安全性、社会的影響を設計段階から考える必要がある。
\subsection{2.2 問題解決活動としての設計}
設計は問題解決活動である。しかし、設計問題は開かれており、曖昧に定義され、複数の解き方があることが多い。このような問題は邪悪な問題とも呼ばれる。ここでの邪悪とは悪意があるという意味ではなく、問題を完全に定義するには、一度解こうとしてみる必要があるという意味である。
ソフトウェア開発では、試作や小さな実装を通じて初めて要求が明確になることがある。したがって、設計は最初から完成した正解を出す活動ではなく、理解と解決を往復しながら進む活動である。
\todobox{邪悪な問題としての設計を示す図を入れる。}{白背景、モノクロ、A4本文幅の概念図を作成する。左に「曖昧な問題」、中央に「試しに解く」、右に「問題理解が深まる」、さらに「よりよい解を作る」へ循環する矢印を描く。「一度解くことで問題が定義される」という注記を入れる。日本語、講義ノート風。}
\section{3 抽象化とカプセル化}
抽象化は、問題解決に不可欠な技法である。対象のすべての細部を同時に扱うのではなく、目的に関係する情報だけに注目する。抽象化は曖昧にすることではない。むしろ、どの水準で正確に議論するかを決めることである。
\subsection{3.1 抽象化の水準}
抽象化では、全体像の一つの水準に集中する。たとえば、Webサービスを考えるとき、利用者体験の水準、業務処理の水準、APIの水準、データベースの水準、物理サーバの水準は異なる。今どの水準で議論しているかを明確にしないと、議論が混乱する。
抽象化の水準は、現実の部品と一対一に対応するとは限らない。むしろ、標準化されたインタフェースを境界として設けることが重要である。APIのような標準インタフェースは、移植性、統合の容易さ、再利用性を高める。
\subsection{3.2 カプセル化}
カプセル化は、抽象化を実現する仕組みである。ある水準で作業するとき、その上下の水準の詳細を隠す。利用者は、対象の内部表現や処理手順を知らなくても、提供された操作を通じて対象を使える。
オブジェクト指向で、オブジェクトの内部データを隠し、公開されたメソッドだけを通じて操作するのは典型例である。カプセル化により、内部実装を変更しても、外部インタフェースが保たれていれば利用側への影響を抑えられる。
\subsection{3.3 階層}
抽象化の水準は、しばしば階層として整理される。階層は、逐次的な一列の構造、木構造、多対多の構造を取ることがある。ただし、階層には循環があってはならない。循環すると、どの水準がどの水準を支えるかが不明確になる。
タスク分解も階層を作る。組織の大きな目標を小さなタスクへ分け、それぞれをさらにサブタスクへ分けると、作業の責任範囲と依存関係が理解しやすくなる。
\subsection{3.4 代替的な抽象化}
同じ問題に対して、複数の抽象化を同時に持つことが有用な場合がある。たとえば同じソフトウェアについて、クラス図、状態遷移図、シーケンス図を使うことがある。これらは上下関係ではなく、異なる視点から同じ対象を説明する図である。
代替的な抽象化は理解を助ける一方で、同期が難しい。クラス図で変更した内容が状態遷移図やシーケンス図にも反映されているかを確認しなければならない。
\todobox{抽象化の水準と代替的な抽象化の図を入れる。}{白背景、モノクロ、A4本文幅の図を作成する。左側に階層として「利用者体験」「業務処理」「API」「データ」「実行環境」を縦に配置する。右側に同じ水準を表す「クラス図」「状態遷移図」「シーケンス図」を横並びで配置し、これらは階層ではなく相補的な視点であると注記する。日本語、細線、講義ノート風。}
\section{4 経験的方法と実験技法}
工学では、候補解やそのモデルを提案し、それを観察・実験・試験によって調べる。経験的方法は、観測値のばらつきを理解し、ばらつきの源を特定し、よりよい判断を行うための方法である。
工学でよく使われる経験的研究には、計画実験、観察研究、回顧研究がある。これらは、どれが常に優れているという関係ではない。操作できる条件、必要な証拠、利用可能なデータ、現場の制約に応じて使い分ける。
\subsection{4.1 計画実験}
計画実験は、明確な仮説を立て、独立変数を操作し、従属変数への効果を測る方法である。独立変数とは研究者が操作する要因であり、従属変数とはその結果として観察する値である。
たとえば、新しい静的解析ツールを使うと欠陥検出率が上がるかを調べる場合、ツール使用の有無が独立変数、検出された欠陥数やレビュー時間が従属変数になる。独立変数の値の組合せを処置と呼ぶ。
\subsection{4.2 観察研究}
観察研究、または事例研究は、現実の文脈の中で過程や現象を観察する方法である。実験が文脈を意図的に切り離すのに対して、観察研究は文脈を含めて理解する。
観察研究は、なぜ、どのようにという問いに向いている。また、参加者の行動を操作できない場合、現象と文脈の境界がはっきりしない場合に有用である。
\subsection{4.3 回顧研究}
回顧研究は、過去に蓄積されたデータを分析する方法である。履歴データから、変数間の関係、将来の予測、傾向を見つける。ソフトウェアでは、障害履歴、変更履歴、テスト結果、レビュー記録などが対象になり得る。
ただし、分析の品質は過去データの品質に依存する。過去データが欠けている、測り方が途中で変わっている、誤って記録されている場合、分析結果にも限界がある。
\begin{longtable}{>{\raggedright\arraybackslash}p{0.17\linewidth}>{\raggedright\arraybackslash}p{0.23\linewidth}>{\raggedright\arraybackslash}p{0.29\linewidth}>{\raggedright\arraybackslash}p{0.24\linewidth}}
\toprule
\textbf{研究の種類} & \textbf{主な特徴} & \textbf{ソフトウェアでの例} & \textbf{注意点} \\
\midrule
\endfirsthead
\toprule
\textbf{研究の種類} & \textbf{主な特徴} & \textbf{ソフトウェアでの例} & \textbf{注意点} \\
\midrule
\endhead
計画実験 & 仮説を立て、条件を操作して効果を見る。 & ツール使用の有無でレビュー効率を比較する。 & 明確な仮説と実験設計が必要である。 \\
観察研究 & 現実の文脈を含めて観察する。 & アジャイルチームの振り返り運用を観察する。 & 因果関係の断定には注意が必要である。 \\
回顧研究 & 過去データを分析する。 & 障害履歴から再発しやすい欠陥種類を探す。 & データの欠落や測定方法の変化に注意する。 \\
\bottomrule
\end{longtable}
\todobox{経験的研究の三分類を比較する図を入れる。}{白背景、モノクロ、A4本文幅の比較図を作成する。三つの箱「計画実験」「観察研究」「回顧研究」を横並びにし、それぞれに「条件を操作する」「現場を観察する」「過去データを見る」と短く書く。下にソフトウェア例を一つずつ示す。日本語、講義ノート風。}
\section{5 統計分析}
工学者は、製品や過程の特性がどのようにばらつくかを理解する必要がある。多くの調査では標本を使うが、知りたいのは母集団である。したがって、信頼できるデータを集め、標本から母集団について妥当な結論を出すために統計分析が必要である。
\subsection{5.1 分析単位、母集団、標本}
\begin{longtable}{>{\raggedright\arraybackslash}p{0.21\linewidth}>{\raggedright\arraybackslash}p{0.40\linewidth}>{\raggedright\arraybackslash}p{0.31\linewidth}}
\toprule
\textbf{概念} & \textbf{説明} & \textbf{ソフトウェアでの例} \\
\midrule
\endfirsthead
\toprule
\textbf{概念} & \textbf{説明} & \textbf{ソフトウェアでの例} \\
\midrule
\endhead
分析単位 & 観察や測定の対象になる単位である。 & 利用者、機能、欠陥、コミット、チケットなど。 \\
母集団 & 関心のある全対象の集合である。 & ある製品を使う全利用者、全障害報告、全ソースファイルなど。 \\
標本 & 母集団から選ばれた一部である。 & 利用者100人の調査、過去3か月の障害報告など。 \\
確率標本抽出 & 各単位が選ばれる確率を事前に定義して標本を取る方法である。 & 全利用者から無作為に調査対象を選ぶ。 \\
研究母集団と対象母集団 & 実際に調べた集団と、結果を一般化したい集団は同じとは限らない。 & 過去データから将来の欠陥傾向を推測する場合。 \\
\bottomrule
\end{longtable}
確率変数とは、観測や測定の結果として値を取る量である。値の集合が有限または数えられる場合は離散型、連続的な値を取る場合は連続型である。確率変数の取り得る値の部分集合を事象という。
\begin{longtable}{>{\raggedright\arraybackslash}p{0.17\linewidth}>{\raggedright\arraybackslash}p{0.46\linewidth}>{\raggedright\arraybackslash}p{0.29\linewidth}}
\toprule
\textbf{分布} & \textbf{何を表すか} & \textbf{例} \\
\midrule
\endfirsthead
\toprule
\textbf{分布} & \textbf{何を表すか} & \textbf{例} \\
\midrule
\endhead
二項分布 & 独立した試行で成功回数を数える。成功確率は一定と仮定する。 & テストケース100件のうち失敗する件数。 \\
ポアソン分布 & 時間や空間における事象の発生回数を表す。 & 1時間あたりの障害発生件数。 \\
正規分布 & 多数の値を取る連続または離散の変数を近似するためによく使う。 & 作業時間や測定誤差の近似。 \\
\bottomrule
\end{longtable}
分布は、平均、標準偏差、発生率、成功確率などの母数によって特徴付けられる。母数は通常未知であるため、標本から推定する。標本平均など、母数を推定するための標本から計算した値を統計量という。
推定には、点推定と区間推定がある。点推定は一つの値で母数を推定する。区間推定は、標本のばらつきを考慮して、母数が入りそうな区間を示す。推定量には、効率性、一致性、偏りの少なさなどの性質が求められる。
仮説検定では、変化がないという帰無仮説 H0 と、関心のある対立仮説 H1 を立てる。標本から統計量を計算し、あらかじめ定めた棄却域に入るかどうかで、H0を棄却するか判断する。
\begin{longtable}{>{\raggedright\arraybackslash}p{0.20\linewidth}>{\raggedright\arraybackslash}p{0.36\linewidth}>{\raggedright\arraybackslash}p{0.36\linewidth}}
\toprule
\textbf{自然の状態} & \textbf{統計的判断：H0を採択} & \textbf{統計的判断：H0を棄却} \\
\midrule
\endfirsthead
\toprule
\textbf{自然の状態} & \textbf{統計的判断：H0を採択} & \textbf{統計的判断：H0を棄却} \\
\midrule
\endhead
H0が真 & 正しい判断 & 第1種の誤り。真のH0を誤って棄却する。 \\
H0が偽 & 第2種の誤り。偽のH0を誤って採択する。 & 正しい判断 \\
\bottomrule
\end{longtable}
検定では、第1種の誤りの確率を一定値、たとえば5\%以下に保ちつつ、第2種の誤りを減らし、検出力を高めることを目指す。
\todobox{母集団と標本の関係図を入れる。}{白背景、モノクロ、A4本文幅の図を作成する。大きな円を「母集団」、その中の小さな点群を「標本」として示す。矢印で「標本から統計量を計算」「母数を推定」「対象母集団へ一般化」と書く。研究母集団と対象母集団が異なる場合の注意も注記する。日本語、細線、講義ノート風。}
\todobox{仮説検定の第1種・第2種の誤り表を図示する。}{白背景、モノクロ、A4本文幅の2×2表を作成する。行は「H0が真」「H0が偽」、列は「H0を採択」「H0を棄却」。セルに「正しい判断」「第1種の誤り」「第2種の誤り」「正しい判断」を置く。下に「第1種の誤りを抑え、検出力を高める」と注記する。日本語、講義ノート風。}
\subsection{5.2 相関と回帰}
相関は、二つの変数の線形関係の強さを表す。相関係数はマイナス1からプラス1までの値を取る。プラスなら一方が増えると他方も増える傾向、マイナスなら一方が増えると他方は減る傾向である。
\infobox{注意}{相関は因果を意味しない。二つの変数に関係があっても、一方が他方を引き起こしたとは限らない。第三の要因が両方に影響している場合もある。}
回帰は、変数間の関係の強さを用いて、一方から他方を推定する分析である。決定係数は0から1までの値を取り、1に近いほど関係が強い。ソフトウェアでは、規模から工数を推定する、欠陥数から修正時間を推定する、といった場面で使われる。
\todobox{相関と回帰の違いを示す散布図を入れる。}{白背景、モノクロ、A4本文幅の学習用図を作成する。左に散布図と「相関：関係の向きと強さ」、右に散布図へ回帰直線を重ねて「回帰：片方からもう片方を推定」と書く。下に「相関は因果を意味しない」と強調する。日本語、講義ノート風。}
\section{6 モデリング・シミュレーション・プロトタイピング}
モデリング、シミュレーション、プロトタイピングはいずれも抽象化の技法である。モデルはシステムの一部の側面を表す。シミュレーションはモデルを使って実験する。プロトタイプは関心のある側面を具体化した部分的な実体である。
\subsection{6.1 モデリング}
モデルは、現実または想像上の対象の抽象表現である。物理的な模型もモデルであり、CAD図面、数式、UML図、業務フロー図もモデルである。モデルは、何を知っているか、何をまだ知らないかを明らかにする。
\begin{longtable}{>{\raggedright\arraybackslash}p{0.17\linewidth}>{\raggedright\arraybackslash}p{0.48\linewidth}>{\raggedright\arraybackslash}p{0.28\linewidth}}
\toprule
\textbf{モデルの種類} & \textbf{説明} & \textbf{例} \\
\midrule
\endfirsthead
\toprule
\textbf{モデルの種類} & \textbf{説明} & \textbf{例} \\
\midrule
\endhead
類像モデル & 見た目が対象に似ているが、不完全な表現である。 & 地図、地球儀、橋の縮尺模型。 \\
類推モデル & 見た目は似ていなくても、機能的に似た振る舞いをする。 & 風洞実験用の模型飛行機、製造過程のシミュレーション。 \\
記号モデル & 数式や記号を使って対象を高い抽象度で表す。 & 運動方程式、確率モデル、性能モデル。 \\
\bottomrule
\end{longtable}
\subsection{6.2 シミュレーション}
シミュレーションは、現実を単純化したモデルを使って実験する方法である。どの情報を抽象化し、どの情報を残すかを誤ると、シミュレーション結果は無効になる。
シミュレーションには連続シミュレーションと離散シミュレーションがある。ソフトウェア工学では、イベント、キュー、資源、状態、時刻を扱う離散シミュレーションが特に重要である。
離散シミュレーションでは、初期化が大きな問題になる。たとえばキューを空で開始するか、実運用に近い混雑状態から開始するかで結果が変わる。初期条件の選び方は、分析結果に影響する。
\subsection{6.3 プロトタイピング}
プロトタイピングは、設計中にシステムの初期版や部分版を作り、実現可能性や使い勝手を調べる方法である。要求獲得、利用者インタフェースの設計、機能要求の妥当性確認に使える。
ソフトウェアのプロトタイプは、完成品と同じ品質を持つとは限らない。性能、アーキテクチャ、保守性、安全性は、完成版の水準に達していない場合がある。したがって、何を確認するためのプロトタイプかを明確にし、計画し、監視し、管理する必要がある。
\todobox{モデル・シミュレーション・プロトタイプの関係図を入れる。}{白背景、モノクロ、A4本文幅の概念図を作成する。中央に「対象システム」を置き、左に「モデル：側面を抽象化」、下に「シミュレーション：モデルで実験」、右に「プロトタイプ：部分的に作る」を配置する。三者が対象理解を助けることを矢印で示す。日本語、講義ノート風。}
\section{7 測定}
測定では、何を測るか、どのように測るか、測定結果で何ができるか、なぜ測るかを理解する必要がある。測定は、抽象的な概念から始まり、操作的定義、測定方法、測定結果へ進む。どの段階も曖昧だと、使えるデータにならない。
測定理論では、測定とは、対象の属性に値や記号を割り当てる活動である。ソフトウェアでは、規模、複雑さ、欠陥、信頼性、保守性などの属性が物理量ほど直感的でないため、操作的定義が特に重要である。
\infobox{操作的定義}{操作的定義とは、測定を実際に行う手順を明確にした定義である。たとえば「身長」を測るだけでも、靴を履くか、髪をどう扱うか、いつ測るか、どの精度で読むかを決めなければ値がばらつく。ソフトウェアの「規模」や「複雑さ」では、さらに明確な手順が必要である。}
\subsection{7.1 測定水準（尺度）}
測定値は、名義尺度、順序尺度、間隔尺度、比率尺度のいずれかとして理解できる。尺度によって、許される比較や演算が違う。数値で表されているからといって、足し算、平均、割り算をしてよいとは限らない。
\begin{longtable}{>{\raggedright\arraybackslash}p{0.14\linewidth}>{\raggedright\arraybackslash}p{0.26\linewidth}>{\raggedright\arraybackslash}p{0.32\linewidth}>{\raggedright\arraybackslash}p{0.20\linewidth}}
\toprule
\textbf{尺度} & \textbf{意味} & \textbf{許される主な操作} & \textbf{例} \\
\midrule
\endfirsthead
\toprule
\textbf{尺度} & \textbf{意味} & \textbf{許される主な操作} & \textbf{例} \\
\midrule
\endhead
名義尺度 & 分類名を付けるだけの尺度である。順序はない。 & 同じか違うかを判定する。同じ分類の数を数える。 & 職名、車種、開発ライフサイクルの種類。 \\
順序尺度 & 分類に順序がある尺度である。差の大きさは意味を持たない。 & 同じか違うか、大小比較、中央値の計算。加減乗除や平均は不適切である。 & 順位、重大度、同意度、CMMI成熟度段階。 \\
間隔尺度 & 隣り合う値の差が一定の尺度である。ゼロは任意である。 & 大小比較、加減、平均、標準偏差、相関、回帰。乗除は意味を持たない。 & 摂氏・華氏温度、暦日、北米式靴サイズ。 \\
比率尺度 & ゼロが属性の不在を表す尺度である。 & すべての算術演算と統計処理が可能である。 & ケルビン温度、金額、分岐構造の数。 \\
絶対尺度 & 変換の余地がない比率尺度である。 & 数そのものを扱う。 & プロジェクトに参加する技術者数。 \\
\bottomrule
\end{longtable}
\todobox{測定尺度の階段図を入れる。}{白背景、モノクロ、A4本文幅の階段図を作成する。下から「名義尺度」「順序尺度」「間隔尺度」「比率尺度」「絶対尺度」と積み上げる。各段に「分類」「順序」「差」「真のゼロ」「唯一の数え上げ」と短く書く。横に「上に行くほど許される演算が増える」と注記する。日本語、講義ノート風。}
\subsection{7.2 測定理論がプログラミング言語に与える含意}
一般的なプログラミング言語は、整数、浮動小数点数、文字、文字列、列挙型などを提供する。しかし、多くの言語は測定理論を理解していない。整数型や浮動小数点型は、加減乗除や比較を自由に許す。
たとえば、CMMI成熟度段階を1から5の整数で表すと、言語は平均や掛け算を許してしまう。しかし成熟度段階は順序尺度であり、段階4が段階2の二倍よいとは言えない。同様に、車種名を文字列で表すと辞書順比較は可能だが、その順序に意味があるとは限らない。
将来の言語は、名義尺度、順序尺度、間隔尺度、比率尺度の意味を型として扱い、不適切な演算を警告または禁止するべきである。それまでは、技術者が測定尺度を理解し、コードレビューで不適切な操作を見つける必要がある。
\todobox{測定尺度とプログラミング型のずれを示す図を入れる。}{白背景、モノクロ、A4本文幅の図を作成する。左に「測定理論：名義・順序・間隔・比率」、右に「一般的な型：int・double・string・enum」を配置する。中央に「型は演算を許すが、尺度は意味を制約する」と書く。CMMI段階の平均計算が危険な例を小さく添える。日本語、講義ノート風。}
\subsection{7.3 直接測度と導出測度}
直接測度は、対象から直接得る測度である。たとえば、見つかった欠陥数、コミット数、レビュー件数などである。導出測度は、複数の直接測度を組み合わせて得る。たとえば、欠陥1件あたりの平均修正時間、千行あたりの欠陥密度などである。
導出測度では、元になった尺度の制約を守る必要がある。異なる尺度を組み合わせるとき、結果の尺度は最も低い尺度の制約を超えられない。高い尺度として扱いたければ、追加の定義や測定努力が必要である。
\begin{longtable}{>{\raggedright\arraybackslash}p{0.18\linewidth}>{\raggedright\arraybackslash}p{0.41\linewidth}>{\raggedright\arraybackslash}p{0.32\linewidth}}
\toprule
\textbf{測度} & \textbf{説明} & \textbf{例} \\
\midrule
\endfirsthead
\toprule
\textbf{測度} & \textbf{説明} & \textbf{例} \\
\midrule
\endhead
直接測度 & 対象から直接数える、読む、測る。 & 欠陥数、処理時間、レビュー件数。 \\
導出測度 & 複数の測度を組み合わせて作る。 & 欠陥密度、平均修正時間、稼働率。 \\
\bottomrule
\end{longtable}
\subsection{7.4 信頼性と妥当性}
測定方法を考えるときは、その方法がよい品質で概念を測っているかを問う必要がある。そのための基本概念が信頼性と妥当性である。
\begin{longtable}{>{\raggedright\arraybackslash}p{0.17\linewidth}>{\raggedright\arraybackslash}p{0.48\linewidth}>{\raggedright\arraybackslash}p{0.28\linewidth}}
\toprule
\textbf{概念} & \textbf{説明} & \textbf{確認の観点} \\
\midrule
\endfirsthead
\toprule
\textbf{概念} & \textbf{説明} & \textbf{確認の観点} \\
\midrule
\endhead
信頼性 & 同じ対象を何度測っても一貫した結果が得られる程度である。 & 測定値のばらつきが小さいか。変動係数が小さいか。 \\
妥当性 & 測りたい概念を本当に測っている程度である。 & 構成概念妥当性、基準妥当性、内容妥当性から見る。 \\
\bottomrule
\end{longtable}
\subsection{7.5 信頼性の評価}
信頼性の評価方法には、再検査法、代替形式法、折半法、内的一貫性法がある。最も単純なのは再検査法である。同じ対象に同じ測定方法を二回適用し、二回の結果の相関係数から測定方法の信頼性を評価する。
\subsection{7.6 目標・質問・測度パラダイム：なぜ測るのか}
目標・質問・測度パラダイム、すなわちGQMは、測定は意思決定を支えるために行うべきだという考え方である。まず目標を決め、その目標に答える質問を決め、質問に答えるための測度を決める。
現実の組織では、測りやすいから、グラフにすると面白いからという理由だけで測定が集められることがある。これは単なる好奇心のための測定であり、意思決定に使われないなら時間と労力の浪費である。
\todobox{GQMの流れ図を入れる。}{白背景、モノクロ、A4本文幅の横長図を作成する。左から「目標：何を改善したいか」「質問：何を知れば判断できるか」「測度：何を測れば答えられるか」「意思決定：どう行動するか」を矢印でつなぐ。下に「測りやすいものではなく、判断に必要なものを測る」と注記する。日本語、講義ノート風。}
\section{8 標準}
標準は、比較のための基準、許容範囲を示す分類、要求される水準などを定めるものである。工学では、製品、過程、材料が受け入れ可能な品質を持つように、要求、仕様、指針を提供する。
標準への適合は、組織の製品や過程が一定の要求を満たすことを社会に示す手段になる。ただし、標準が有用であるためには、適合することが実際または知覚上の価値を生む必要がある。
標準は、購入者を保護し、事業者を保護し、ソフトウェア工学で使う方法や手順を明確にし、共通用語と期待を提供する。国際的な標準化組織には、ITU、IEC、IEEE、ISOなどがある。国や地域にも標準化組織がある。
標準は、多くの場合、合意形成に基づいて作られる。技術者は、自分の領域で適用される標準を把握し、標準の変更にも追随する必要がある。法律や契約が特定標準への適合を求める場合、標準は設計制約の最小要求になる。
\todobox{標準と適合の概念図を入れる。}{白背景、モノクロ、A4本文幅の概念図を作成する。中央に「標準」を置き、そこから「要求」「仕様」「指針」「共通用語」「品質水準」へ矢印を出す。右側に「適合している製品」と「適合していない製品」を分ける境界線を描く。下に「標準は設計制約になり得る」と注記する。日本語、講義ノート風。}
\section{9 根本原因分析}
根本原因分析は、望ましくない結果の背後にある原因を特定する問題解決方法群である。症状だけを直すのではなく、なぜ、どのようにその結果が起きたかを理解し、再発を防ぐ行動につなげる。
ソフトウェアプロジェクトでは、根本原因分析は多くの場面で役立つ。工学活動で本当に解くべき問題を見つける。リスクの背後にある要因を明らかにする。プロセス改善の機会を見つける。繰り返し発生する欠陥の原因を調べる。
\subsection{9.1 根本原因分析技法}
\begin{longtable}{>{\raggedright\arraybackslash}p{0.20\linewidth}>{\raggedright\arraybackslash}p{0.51\linewidth}>{\raggedright\arraybackslash}p{0.21\linewidth}}
\toprule
\textbf{技法} & \textbf{説明} & \textbf{使いどころ} \\
\midrule
\endfirsthead
\toprule
\textbf{技法} & \textbf{説明} & \textbf{使いどころ} \\
\midrule
\endhead
変更分析 & うまくいった状況とうまくいかなかった状況を比較し、差分に原因を探す。 & リリース後に突然障害が増えた場合。 \\
5つのなぜ & 望ましくない結果から始め、なぜを繰り返して深い原因に近づく。 & 単純な障害やプロセス問題の初期分析。 \\
原因結果図 & 特性要因図、石川図、魚骨図とも呼ばれる。原因を人、過程、道具、材料、測定、環境などに分けて整理する。 & 原因候補を広く出す場合。 \\
故障木解析 & 原因と結果のAND/OR関係を明示する、より形式的な原因分析である。 & 安全性や信頼性の分析。 \\
故障モード影響解析 & 故障し得る要素から出発し、その影響がどう連鎖するかを前向きに調べる。 & 設計段階の予防的な分析。 \\
原因マップ & 証拠に基づいて、原因から結果、さらに組織目標への影響までを構造化する。 & 厳密な再発防止分析。 \\
現実ツリー & 論理規則に従って、現状の望ましくない結果と原因を木として表す。 & 複数の問題が絡む改善課題。 \\
人間性能評価 & 入力検出、入力理解、行動選択、行動実行のどこで失敗したかを見る。 & 運用ミス、認知負荷、疲労、思い込みが関わる問題。 \\
\bottomrule
\end{longtable}
\todobox{根本原因分析技法の比較図を入れる。}{白背景、モノクロ、A4本文幅の比較図を作成する。左から「5つのなぜ」「魚骨図」「故障木解析」「FMEA」「原因マップ」を並べる。上に「後ろ向きに原因を探す」から「前向きに影響を追う」までの軸を置く。各技法の特徴を短く書く。日本語、講義ノート風。}
\subsection{9.2 根本原因に基づく改善}
根本原因を特定しても、何も改善しなければ意味がない。根本原因分析は、より大きなプロセス改善の一部として使う必要がある。
\begin{longtable}{>{\raggedright\arraybackslash}p{0.23\linewidth}>{\raggedright\arraybackslash}p{0.69\linewidth}}
\toprule
\textbf{段階} & \textbf{内容} \\
\midrule
\endfirsthead
\toprule
\textbf{段階} & \textbf{内容} \\
\midrule
\endhead
解く問題を選ぶ & パレート分析、頻度と重大度、統計的工程管理などを使い、優先度の高い望ましくない結果を選ぶ。 \\
証拠を集める & 証言、過程、標準、仕様、報告、履歴傾向、実験、試験などを集める。 \\
根本原因を特定する & 一つまたは複数の根本原因分析技法を使う。 \\
是正処置を選ぶ & 再発防止、組織が制御可能、目標に合う、別問題を起こさない、費用対効果が高いという条件で選ぶ。 \\
是正処置を実施する & 選んだ対策を実際に導入する。 \\
効果を観察する & 効率性と有効性を確認し、必要なら改善する。 \\
\bottomrule
\end{longtable}
\todobox{根本原因に基づく改善サイクルを入れる。}{白背景、モノクロ、A4本文幅の循環図を作成する。「問題を選ぶ」「証拠を集める」「根本原因を特定する」「是正処置を選ぶ」「実施する」「効果を観察する」を円環に配置する。中央に「再発防止」と書く。日本語、講義ノート風。}
\section{10 インダストリー4.0とソフトウェア工学}
インダストリー4.0は、製造業を大きく変える概念である。特に、AIを活用した個別化された製造、デジタル化、標準システムとの統合が重要になる。これにより、費用、品質、効率に便益が期待される。
インダストリー4.0では、ソフトウェアが中心的な構成要素になる。多くの機器が相互接続され、データを送り、命令やデータを受け取り、さらに動作する。したがって、堅牢で知的なシステムを作るには、ソフトウェアを工学的に作ることが不可欠である。
継続的ソフトウェア工学は、継続的な計画、アーキテクチャ設計、開発、統合、配備、レビュー・改訂を支える。製造が継続的に変化するなら、ソフトウェア開発も継続的に学習し、更新する必要がある。
\begin{longtable}{>{\raggedright\arraybackslash}p{0.28\linewidth}>{\raggedright\arraybackslash}p{0.64\linewidth}}
\toprule
\textbf{技術} & \textbf{ソフトウェア工学での意味} \\
\midrule
\endfirsthead
\toprule
\textbf{技術} & \textbf{ソフトウェア工学での意味} \\
\midrule
\endhead
IoT & 機器からデータを集め、機器へ命令を返すソフトウェアが必要になる。 \\
ビッグデータ分析 & 大量データを処理し、予測や最適化に使う。 \\
AI・機械学習 & 予測、検査、自動判断、個別化された製造に使う。 \\
サイバーセキュリティ & 接続機器が増えるほど攻撃面も増えるため、設計から保護が必要である。 \\
クラウドコンピューティング & 計算資源、保存、監視、分析を柔軟に提供する。 \\
複数プラットフォーム向けアプリ & 工場、モバイル、クラウド、組込み機器をまたいだ利用者体験を支える。 \\
量子コンピューティング & 複雑な計算をより高速かつ費用対効果高く実行する可能性がある。 \\
\bottomrule
\end{longtable}
今後のソフトウェアは、自己学習的で先回りする性質を強める可能性がある。利用者が何を望むかを予測し、状況に応じて振る舞いを変えるためには、測定、実験、モデル、標準、根本原因分析のすべてが重要になる。
\todobox{インダストリー4.0とソフトウェア工学の関係図を入れる。}{白背景、モノクロ、A4本文幅の概念図を作成する。中央に「ソフトウェア工学」を置き、周囲に「IoT」「AI・機械学習」「ビッグデータ」「クラウド」「サイバーセキュリティ」「複数プラットフォーム」「継続的工学」を配置する。外側に「製造」「運用」「保守」「改善」を円環で示す。日本語、講義ノート風。}
\section{11 トピックと参考文献の対応}
この表は、SWEBOKが示すトピック対参考資料の行列を、学習用に簡略化したものである。詳しく学ぶときは、各トピックに対応する文献を読む。
\begin{longtable}{>{\raggedright\arraybackslash}p{0.33\linewidth}>{\raggedright\arraybackslash}p{0.59\linewidth}}
\toprule
\textbf{トピック} & \textbf{主な参考資料} \\
\midrule
\endfirsthead
\toprule
\textbf{トピック} & \textbf{主な参考資料} \\
\midrule
\endhead
1 工学プロセス & Tockey, Return on Software \\
2 工学設計 & Voland, Engineering by Design; McConnell, Code Complete \\
3 抽象化とカプセル化 & McConnell, Code Complete \\
4 経験的方法と実験技法 & Montgomery and Runger, Applied Statistics and Probability for Engineers \\
5 統計分析 & Montgomery and Runger; Null and Lobur \\
6 モデリング・シミュレーション・プロトタイピング & Voland; Cheney and Kincaid; Sommerville; Fairley \\
7 測定 & Tockey; Voland; Fairley; Kan \\
8 標準 & Voland; Moore \\
9 根本原因分析 & Voland; Kan; Tockey; Ishikawa; Gano; Goldratt \\
10 インダストリー4.0とソフトウェア工学 & Nakagawa et al., Continuous Systems and Software Engineering for Industry 4.0 \\
\bottomrule
\end{longtable}
\section{12 発展的な読み物}
\begin{longtable}{>{\raggedright\arraybackslash}p{0.31\linewidth}>{\raggedright\arraybackslash}p{0.61\linewidth}}
\toprule
\textbf{文献} & \textbf{読むと得られること} \\
\midrule
\endfirsthead
\toprule
\textbf{文献} & \textbf{読むと得られること} \\
\midrule
\endhead
Abran, Software Metrics and Software Metrology & 測度、測定方法、測定結果という用語を厳密に使い分け、ソフトウェア測定をより正確に理解できる。測定の節を深めるのに有用である。 \\
Vincenti, What Engineers Know and How They Know It & 実際の工学事例を通じて、工学者が何を知り、どのように知識を得るかを理解できる。ソフトウェア以外の工学と比較するのに有用である。 \\
\bottomrule
\end{longtable}
\section{13 参考文献}
\begin{enumerate}[leftmargin=2em]
\item ISO/IEC/IEEE 24765:2017, Systems and Software Engineering – Vocabulary.
\item S. Tockey, Return on Software: Maximizing the Return on Your Software Investment, Addison-Wesley, 2004.
\item G. Voland, Engineering by Design, 2nd ed., Prentice Hall, 2003.
\item Canadian Engineering Accreditation Board, 2021 Accreditation Criteria and Procedures, Engineers Canada, 2021.
\item ABET, Criteria for Accrediting Engineering Programs, 2022–2023, 2021.
\item S. McConnell, Code Complete, 2nd ed., Microsoft Press, 2004.
\item E. W. Dijkstra, The Humble Programmer, Communications of the ACM, 1972.
\item D. C. Montgomery and G. C. Runger, Applied Statistics and Probability for Engineers, 7th ed., Wiley, 2018.
\item L. Null and J. Lobur, The Essentials of Computer Organization and Architecture, 5th ed., Jones and Bartlett Publishers, 2018.
\item E. W. Cheney and D. R. Kincaid, Numerical Mathematics and Computing, 7th ed., Brooks/Cole, 2020.
\item I. Sommerville, Software Engineering, 10th ed., Addison-Wesley, 2016.
\item R. E. Fairley, Managing and Leading Software Projects, Wiley-IEEE Computer Society Press, 2009.
\item S. H. Kan, Metrics and Models in Software Quality Engineering, 2nd ed., Addison-Wesley, 2002.
\item J. W. Moore, The Road Map to Software Engineering: A Standards-Based Guide, Wiley-IEEE Computer Society Press, 2006.
\item K. Ishikawa, Introduction to Quality Control, Productivity Press, 1990.
\item D. Gano, Apollo Root Cause Analysis, 3rd ed., Apollonian Publications, 2007.
\item E. Goldratt, It’s Not Luck, North River Press, 1994.
\item A. Abran, Software Metrics and Software Metrology, Wiley-IEEE Computer Society Press, 2010.
\item W. G. Vincenti, What Engineers Know and How They Know It, Johns Hopkins University Press, 1993.
\item E. Y. Nakagawa, P. O. Antonio, F. Schnicke, T. Kuhn, and P. Liggesmeyer, Continuous Systems and Software Engineering for Industry 4.0: A disruptive view, Information and Software Technology, 2021.
\end{enumerate}
\section{14 画像 TODO 一覧}
本文に配置した画像TODOを再掲する。実際に図を作る場合は、各TODOのプロンプトをそのまま画像生成に使うと、A4資料に合う図を作りやすい。
\begin{enumerate}[leftmargin=2em]
\item 第18章の全体地図を入れる。
\item 工学プロセスの流れ図を入れる。
\item 邪悪な問題としての設計を示す図を入れる。
\item 抽象化の水準と代替的な抽象化の図を入れる。
\item 経験的研究の三分類を比較する図を入れる。
\item 母集団と標本の関係図を入れる。
\item 仮説検定の第1種・第2種の誤り表を図示する。
\item 相関と回帰の違いを示す散布図を入れる。
\item モデル・シミュレーション・プロトタイプの関係図を入れる。
\item 測定尺度の階段図を入れる。
\item 測定尺度とプログラミング型のずれを示す図を入れる。
\item GQMの流れ図を入れる。
\item 標準と適合の概念図を入れる。
\item 根本原因分析技法の比較図を入れる。
\item 根本原因に基づく改善サイクルを入れる。
\item インダストリー4.0とソフトウェア工学の関係図を入れる。
\end{enumerate}
\section{15 章末確認問題}
\begin{enumerate}[leftmargin=2em]
\item 工学プロセスの六つの段階を、自分の言葉で説明せよ。
\item 工学設計が一つの正解を探す活動ではない理由を説明せよ。
\item 抽象化とカプセル化の違いを、APIを例に説明せよ。
\item 計画実験、観察研究、回顧研究はどのように使い分けるか。
\item 分析単位、母集団、標本の違いを、利用者満足度調査を例に説明せよ。
\item 相関が因果を意味しない理由を述べよ。
\item 名義尺度、順序尺度、間隔尺度、比率尺度で許される演算の違いを説明せよ。
\item CMMI成熟度段階を整数で表したときに起こり得る測定理論上の問題を説明せよ。
\item GQMが「測りやすいものを測る」発想と異なる理由を説明せよ。
\item 標準が設計制約になる場合の例を挙げよ。
\item 根本原因分析と表面的な対症療法の違いを説明せよ。
\item インダストリー4.0でソフトウェア工学が重要になる理由を説明せよ。
\end{enumerate}
\infobox{解答の方向性}{1は真の問題、基準、候補解、評価、選択、監視を含める。2は制約下で複数解を比較する点を述べる。3は細部を隠して水準を定める抽象化と、インタフェース越しに内部を隠すカプセル化を区別する。7と8では、数値であっても尺度が違えば平均や乗除が不適切になる点を説明する。10から12では、標準、再発防止、接続された製造システムという語を使うとよい。}
\end{document}